\ifdefined\gkver\else\def\gkver{student}\fi
\documentclass[\gkver]{gaokaozhenti}
\begin{document}

\chapter{2004年上海}

\begin{enumerate}
\item
%% number: 1
%% paperName: 2004年上海高考第 1 题 5 分
%% typeId: 2
%% type: 多选题
%% chapter: 光学
%% point: 光的电磁说
%% method: 无
%% score: 5
%% degree: 850
%% duplicateId: 0
%% body:
下列说法中正确的是\xzanswer{AD}
\fourchoices[answer=AD]
{光的干涉和衍射现象说明光具有波动性}
{光的频率越大，波长越大}
{光的波长越大，光子的能量越大}
{光在真空中的传播速度为 $3.00\times10^{8}\Ums$}

%% answer: AD
%% memo:
\memoanswer{
光的干涉和衍射是特有的现象，因此光具有波动性，故 A 正确；由公式 $\lambda=\frac{v}{f}$ 可知，光的波长由频率和波速决定，故 B 错误；由 $E=h\nu$ 知，光子的能量由频率决定，故 C 错误；任何光在真空中的传播速度均为 $3\times10^{8}\Ums$，故 D 正确。
}

\item
%% number: 2
%% paperName: 2004年上海高考第 2 题 5 分
%% typeId: 2
%% type: 多选题
%% chapter: 其他
%% point: 物理学史
%% method: 无
%% score: 5
%% degree: 850
%% duplicateId: 0
%% body:
下列说法中正确的是\xzanswer{CD}
\fourchoices[answer=CD]
{玛丽$\cdot$居里首先提出原子的核式结构}
{卢瑟福在 $\alpha$ 粒子散射实验中发现了电子}
{查德威克在原子核人工转变的实验中发现了中子}
{爱因斯坦为解释光电效应的实验规律提出了光子说}

%% answer: CD
%% memo:
\memoanswer{
原子的核式结构学说是卢瑟福首先提出的，故 A 错误；电子是汤姆生发现的，故 B 错误；查德威克用高能粒子轰击核，发现了中子，故 C 正确；爱因斯坦提出了光子说，从而解释了光电效应，故 D 正确。
}

\item
%% number: 3
%% paperName: 2004年上海高考第 3 题 5 分
%% typeId: 2
%% type: 多选题
%% chapter: 曲线运动
%% point: 人造地球卫星
%% method: 无
%% score: 5
%% degree: 650
%% duplicateId: 0
%% body:
火星有两颗卫星，分别为火卫一和火卫二，它们的轨道近似为圆，已知火卫一的周期为 7 小时 39 分，火卫二的周期为 30 小时 18 分，则两颗卫星相比\xzanswer{AC}
\fourchoices[answer=AC]
{火卫一距火星表面较近}
{火卫二的角速度较大}
{火卫一的运动速度较大}
{火卫二的向心加速度较大}

%% answer: AC
%% memo:
\memoanswer{
卫星做匀速圆周运动，由 $G\frac{Mm}{r^{2}}=m\frac{4\pi^{2}}{T^{2}}r$ 得 $r=\sqrt[3]{\frac{GMT^{2}}{4\pi^{2}}}$，因 $T_{1}<T_{2}$，故 $r_{1}<r_{2}$，故 A 正确；由 $\omega=\frac{2\pi}{T}$ 可知，$\omega_{1}>\omega_{2}$，故 B 错误；由 $G\frac{Mm}{r^{2}}=m\frac{v^{2}}{r}$ 得 $v=\sqrt{\frac{GM}{r}}$，可知 $v_{1}>v_{2}$，故 C 正确；由 $G\frac{Mm}{r^{2}}=ma$ 知 $a=\frac{GM}{r^{2}}$，可知 $a_{1}>a_{2}$，故 D 错误。
}

\item
%% number: 4
%% paperName: 2004年上海高考第 4 题 5 分
%% typeId: 2
%% type: 多选题
%% chapter: 电磁感应
%% point: 楞次定律
%% method: 无
%% score: 5
%% degree: 650
%% duplicateId: 0
%% body:
两圆环 A、B 置于同一水平面上，其中 A 为均匀带电绝缘环，B 为导体环，当 A 以如图所示的方向绕中心转动的角速度发生变化时，B 中产生如图所示方向的感应电流，则\xzanswer{BC}
\onepicture[w=4.5cm]{\includesvg{figs/04}}
\fourchoices[answer=BC]
{A 可能带正电且转速减小}
{A 可能带正电且转速增大}
{A 可能带负电且转速减小}
{A 可能带负电且转速增大}

%% answer: BC
%% memo:
\memoanswer{
要产生 B 环中所示的电流，感应磁场垂直纸面向外，由楞次定律知 A 环内的磁场应向里增强或向外减弱，由安培定则可知 BC 正确。
}

\item
%% number: 5
%% paperName: 2004年上海高考第 5 题 5 分
%% typeId: 1
%% type: 单选题
%% chapter: 运动和力
%% point: 牛顿第二定律
%% method: 整体与隔离
%% score: 5
%% degree: 650
%% duplicateId: 0
%% body:
物体 B 放在物体 A 上，A、B 的上下表面均与斜面平行（如图），当两者以相同的初速度靠惯性沿光滑固定斜面 C 向上做匀减速运动时\xzanswer{C}
\onepicture[w=5.5cm]{\includesvg{figs/05}}
\fourchoices[answer=C]
{A 受到 B 的摩擦力沿斜面方向向上}
{A 受到 B 的摩擦力沿斜面方向向下}
{A、B 之间的摩擦力为零}
{A、B 之间是否存在摩擦力取决于 A、B 表面的性质}

%% answer: C
%% memo:
\memoanswer{
设 A 所受摩擦力沿斜面向上，对整体有 $(m_{A}+m_{B})g\sin\alpha=(m_{A}+m_{B})a$，得 $a=g\sin\alpha$；对物体 A 有 $m_{A}g\sin\alpha-f=m_{A}a$，解得 $f=0$，故 C 正确。
}

\item
%% number: 6
%% paperName: 2004年上海高考第 6 题 5 分
%% typeId: 2
%% type: 多选题
%% chapter: 电场
%% point: 电势差与电场强度的关系
%% method: 无
%% score: 5
%% degree: 650
%% duplicateId: 0
%% body:
某静电场沿 $x$ 方向的电势分布如图所示，则\xzanswer{AC}
\onepicture[w=6.5cm]{\includesvg{figs/06}}
\fourchoices[answer=AC]
{在 $0\sim x_{1}$ 之间不存在沿 $x$ 方向的电场}
{在 $0\sim x_{1}$ 之间存在着沿 $x$ 方向的匀强电场}
{在 $x_{1}\sim x_{2}$ 之间存在着沿 $x$ 方向的匀强电场}
{在 $x_{1}\sim x_{2}$ 之间存在着沿 $x$ 方向的非匀强电场}

%% answer: AC
%% memo:
\memoanswer{
由图象知 $O\sim x_{1}$ 之间电势不随距离变化，故不存在沿 $x$ 方向的电场，故 A 正确 B 错误；$x_{1}\sim x_{2}$ 之间有电势降落，且为线性变化，其斜率 $k=\frac{\Delta U}{\Delta x}=E$（不变），所以沿 $x$ 方向存在匀强电场，故 C 正确 D 错误。
}

\item
%% number: 7
%% paperName: 2004年上海高考第 7 题 5 分
%% typeId: 2
%% type: 多选题
%% chapter: 电场
%% point: 带电粒子的偏转
%% method: 无
%% score: 5
%% degree: 450
%% duplicateId: 0
%% body:
光滑水平面上有一边长为 $l$ 的正方形区域处在场强为 $E$ 的匀强电场中，电场方向与正方形一边平行，一质量为 $m$，带电量为 $q$ 的小球由某一边的中点，以垂直于该边的初速度 $v_{0}$ 进入该正方形区域。当小球再次运动到该正方形区域的边缘时，具有的动能可能为\xzanswer{ABC}
\fourchoices[answer=ABC]
{0}
{$\frac{1}{2}mv_{0}^{2}+\frac{1}{2}qEl$}
{$\frac{1}{2}mv_{0}^{2}$}
{$\frac{1}{2}mv_{0}^{2}+\frac{2}{3}qEl$}

%% answer: ABC
%% memo:
\memoanswer{
由题意存在以下几种可能情况（设小球带正电）：

（1）$v$ 与 $E$ 在同一条直线上，且电场力做负功，使小球到达另一边缘时动能恰好为零，故 A 正确；或者未到另一边已减速到零，然后反向加速，到出发点时动能为 $\frac{1}{2}mv_{0}^{2}$，故 C 正确；

（2）$v$ 与 $E$ 在同一条直线上，且电场力做正功，小球到达另一边缘时动能为 $\frac{1}{2}mv_{0}^{2}+qEl$；

（3）$v$ 与 $E$ 垂直，小球做类平抛运动，到边缘时动能为 $\frac{1}{2}mv_{0}^{2}+qEx$，其中 $x\leq\frac{l}{2}$，当 $x=\frac{l}{2}$ 时为 B 选项。综上分析可知 ABC 正确 D 错误。
}

\item
%% number: 8
%% paperName: 2004年上海高考第 8 题 5 分
%% typeId: 2
%% type: 多选题
%% chapter: 功和能
%% point: 动能定理
%% method: 无
%% score: 5
%% degree: 450
%% duplicateId: 0
%% body:
滑块以速率 $v_{1}$ 靠惯性沿固定斜面由底端向上运动。当它回到出发点时速率变为 $v_{2}$，且 $v_{2}<v_{1}$，若滑块向上运动的位移中点为 A，取斜面底端重力势能为零，则\xzanswer{BC}
\fourchoices[answer=BC]
{上升时机械能减小，下降时机械能增大}
{上升时机械能减小，下降时机械能也减小}
{上升过程中动能和势能相等的位置在 A 点上方}
{上升过程中动能和势能相等的位置在 A 点下方}

%% answer: BC
%% memo:
\memoanswer{
由 $v_{2}<v_{1}$ 可知，斜面与木块之间有摩擦力，无论上升还是下降时，都有机械能损失，故 A 错误 B 正确；设物体上升的最大高度为 $h$，初动能为 $E_k$，由于摩擦力做负功，故 $E_k>mgh$，A 点处动能为 $\frac{E_k}{2}$，势能为 $\frac{1}{2}mgh$，故 $\frac{E_k}{2}>\frac{1}{2}mgh$，要使动能等于重力势能，应继续上升，即在 A 点上方，故 C 正确 D 错误。
}

\item
%% number: 9
%% paperName: 2004年上海高考第 9 题 4 分
%% typeId: 3
%% type: 填空题
%% chapter: 光学
%% point: 光电效应
%% method: 无
%% score: 4
%% degree: 850
%% duplicateId: 0
%% body:
在光电效应实验中，如果实验仪器及线路完好，当光照射到光电管上时，灵敏电流计中没有电流通过，可能的原因是：\tkanswer[5cm]{入射光波长太大（或反向电压太大）}。

%% answer:
\jdanswer{
入射光波长太大（或反向电压太大）。
}

%% memo:
\memoanswer{
入射光波长太长（或反向电压太大）。
}

\item
%% number: 10
%% paperName: 2004年上海高考第 10 题 4 分
%% typeId: 3
%% type: 填空题
%% chapter: 机械振动
%% point: 等效单摆
%% method: 等效替代
%% score: 4
%% degree: 650
%% duplicateId: 0
%% body:
在光滑水平面上的 O 点系一长为 $l$ 的绝缘细线，线的一端系一质量为 $m$，带电量为 $q$ 的小球。当沿细线方向加上场强为 $E$ 的匀强电场后，小球处于平衡状态。现给小球一垂直于细线的初速度 $v_{0}$，使小球在水平面上开始运动。若 $v_{0}$ 很小，则小球第一次回到平衡位置所需时间为 \tkanswer[2cm]{$\pi\sqrt{\frac{ml}{qE}}$}。
\onepicture[w=6cm, align=right]{\includesvg{figs/10}}

%% answer:
\jdanswer{
$\pi\sqrt{\frac{ml}{qE}}$。
}

%% memo:
\memoanswer{
对小球受力分析，竖直方向受重力和支持力的作用，相互抵消，对运动没有影响；水平方向受到沿绳的拉力和始终沿场强方向的电场力作用。因 $v_{0}$ 很小，可以认为绳的偏转角不超过 $5^\circ$，可以看成简谐运动，这样的运动情况和单摆类似，是类单摆模型。由题意知，利用等效重力场思想，小球受到的电场力 $qE$ 等效成重力，等效重力加速度为 $g'=\frac{qE}{m}$，由单摆周期公式得 $T=2\pi\sqrt{\frac{l}{g}}$，将 $g'$ 代入周期公式，并注意到所求的时间为半个周期得 $t=\pi\sqrt{\frac{ml}{qE}}$。
}

\item
%% number: 11
%% paperName: 2004年上海高考第 11 题 4 分
%% typeId: 3
%% type: 填空题
%% chapter: 原子和原子核
%% point: 原子核式结构
%% method: 无
%% score: 4
%% degree: 650
%% duplicateId: 0
%% body:
利用扫描隧道显微镜（STM）可以得到物质表面原子排列的图象，从而可以研究物质的构成规律，下面的照片是一些晶体材料表面的 STM 图象，通过观察、比较，可以看到这些材料都是由原子在空间排列而构成，具有一定的结构特征。则构成这些材料的原子在物质表面排列的共同特点是：

\threepicture[w=3.2cm]{figs/11a.png}{figs/11b.png}{figs/11c.png}
\begin{enumerate}
	\item
	\tkanswer[6cm]{在确定方向上原子有规律地排列}；
	\item
	\tkanswer[6cm]{在不同方向上原子的排列规律一般不同，原子排列具有一定对称性等}。
\end{enumerate}

%% answer:
\jdanswer{
\begin{enumerate}
	\item
	在确定方向上原子有规律地排列；
	\item
	在不同方向上原子的排列规律一般不同，原子排列具有一定对称性等。
\end{enumerate}
}

%% memo:
\memoanswer{
在确定方向上原子有规律地排列；在不同方向上原子的排列规律一般不同；原子排列具有一定对称性等。
}

\item
%% number: 12
%% paperName: 2004年上海高考第 12 题 4 分
%% typeId: 3
%% type: 填空题
%% chapter: 电路
%% point: 电路中的图象
%% method: 图象法
%% score: 4
%% degree: 450
%% duplicateId: 0
%% body:
两个额定电压为 $220\UV$ 的白炽灯 $L_{1}$ 和 $L_{2}$ 的 $U-I$ 特性曲线如图所示。$L_{2}$ 的额定功率约为 \tkanswer[1.2cm]{99} $\UW$；现将 $L_{1}$ 和 $L_{2}$ 串联后接到 $220\UV$ 的电源上，电源内阻忽略不计，此时 $L_{2}$ 的实际功率约为 \tkanswer[1.4cm]{17.5} $\UW$。
\onepicture[w=7cm, align=right]{\includesvg{figs/12}}

%% answer:
\jdanswer{
$99\UW$，$17.5\UW$。
}

%% memo:
\memoanswer{
由题中提供的 $U-I$ 特性曲线分析知，灯泡 $L_{2}$ 正常发光时，流过灯泡的电流为 $I_{L_{2}}=0.45\UA$，所以灯泡 $L_{2}$ 的额定功率为 $P_{2}=U_{L_{2}}I_{L_{2}}$，代入数据得 $P_{2}=99\UW$；将两灯泡串联后接在 $220\UV$ 的电源上，此时通过两灯泡的电流相等，两灯泡上的电压和为 $220\UV$，从图中的 $U-I$ 特性曲线分析知，当 $I=0.25\UA$ 时，$U_{L_{1}}'=150\UV$，$U_{L_{2}}'=70\UV$，故此时灯泡 $L_{2}$ 的实际功率为 $P_{2}'=U_{L_{2}}'I=0.25\times70\UW=17.5\UW$。
}

\item
%% number: 13
%% paperName: 2004年上海高考第 13 题 4 分
%% typeId: 3
%% type: 填空题
%% chapter: 机械波
%% point: 波的叠加
%% method: 无
%% score: 4
%% degree: 450
%% duplicateId: 0
%% body:
A、B 两波相向而行，在某时刻的波形与位置如图所示，已知波的传播速度为 $v$，图中标尺每格长度为 $l$。在图中画出又经过 $t=\frac{7l}{v}$ 时的波形。
\onepicture[w=10cm]{\includesvg{figs/13}}

%% answer:
\jdanswer{
如图所示。
}

%% memo:
\memoanswer{
如图所示。

\onepicture[w=10cm]{\includesvg{figs/13_an}}
}

\item
%% number: 14
%% paperName: 2004年上海高考第 14 题 5 分
%% typeId: 4
%% type: 实验题
%% chapter: 直线运动
%% point: 打点计时器公式
%% method: 无
%% score: 5
%% degree: 650
%% duplicateId: 0
%% body:
用打点计时器研究物体的自由落体运动，得到如图一段纸带，测得 $AB=7.65\Ucm$，$BC=9.17\Ucm$。已知交流电频率是 $50\UHz$，则打 B 点时物体的瞬时速度为 \tkanswer[1.5cm]{2.10} $\Ums$。如果实验测出的重力加速度值比公认值偏小，可能的原因是 \tkanswer[4.5cm]{下落过程中存在阻力}。
\onepicture[w=13cm, align=right]{\includesvg{figs/14}}

%% answer:
\jdanswer{
$2.10\Ums$，下落过程中存在阻力。
}

%% memo:
\memoanswer{
由匀变速直线运动的规律知平均速度等于中间时刻速度，故 $v_{B}=\frac{x_{AC}}{t_{AC}}=\frac{(7.65+9.17)\Ucm}{4\times0.02\Us}\approx2.10\Ums$，如果实验测出的重力加速度值比公认值偏小，可能的原因是下落过程中存在阻力等。
}

\item
%% number: 15
%% paperName: 2004年上海高考第 15 题 4 分
%% typeId: 4
%% type: 实验题
%% chapter: 电路
%% point: 测电动势与内阻
%% method: 无
%% score: 4
%% degree: 650
%% duplicateId: 0
%% body:
在测定一节干电池（电动势约为 $1.5\UV$，内阻约为 $2\UO$）的电动势和内阻的实验中，变阻器和电压表各有两个供选：A 电压表量程为 $15\UV$，B 电压表量程为 $3\UV$，A 变阻器为（$20\UO$，$3\UA$），B 变阻器为（$500\UO$，$0.2\UA$）。

电压表应该选 \tkanswer[1cm]{B}（填 A 或 B），这是因为 \tkanswer[4.5cm]{A 电压表量程过大，误差较大}；

变阻器应该选 \tkanswer[1cm]{A}（填 A 或 B），这是因为 \tkanswer[4.5cm]{B 变阻器额定电流过小且调节不便}。

%% answer:
\jdanswer{
B，A 电压表量程过大，误差较大；A，B 变阻器额定电流过小且调节不便。
}

%% memo:
\memoanswer{
电压表应选 B：A 电压表量程过大，误差较大；变阻器应选 A：B 变阻器额定电流过小且调节不便。
}

\item
%% number: 16
%% paperName: 2004年上海高考第 16 题 5 分
%% typeId: 4
%% type: 实验题
%% chapter: 光学
%% point: 光的电磁说
%% method: 等效替代
%% score: 5
%% degree: 650
%% duplicateId: 0
%% body:
下图为一测量灯泡发光强度的装置，AB 是一个有刻度的底座，两端可装两个灯泡。中间带一标记线的光度计可在底座上移动，通过观察可以确定两边灯泡在光度计上的照度是否相同。已知照度与灯泡的发光强度成正比、与光度计到灯泡的距离的平方成反比。现有一个发光强度 $I_{0}$ 的灯泡 a 和一个待测灯泡 b。分别置于底座两端（如图）。
\onepicture[w=8cm]{\includesvg{figs/16}}
\begin{enumerate}
	\item
	怎样测定待测灯泡的发光强度？
	\item
	简单叙述一个可以减小实验误差的方法。
\end{enumerate}

%% answer:
\jdanswer{
\begin{enumerate}
	\item
	接通电源，移动光度计使两边的照度相同，测出距离 $r_{1}$ 和 $r_{2}$，即可得待测灯泡的发光强度 $I_{x}=\frac{r_{2}^{2}}{r_{1}^{2}}I_{0}$。
	\item
	多次测量求平均值。
\end{enumerate}
}

%% memo:
\memoanswer{
（1）接通电源，移动光度计使两边的照度相同，测出距离 $r_{1}$ 和 $r_{2}$，即可得待测灯泡的发光强度 $I_{x}=\frac{r_{2}^{2}}{r_{1}^{2}}I_{0}$。

（2）多次测量求平均值。
}

\item
%% number: 17
%% paperName: 2004年上海高考第 17 题 6 分
%% typeId: 4
%% type: 实验题
%% chapter: 气体性质
%% point: 盖吕萨克定律
%% method: 无
%% score: 6
%% degree: 450
%% duplicateId: 0
%% body:
一根长约为 $30\Ucm$、管内截面积为 $S=5.0\times10^{-6}\Umq$ 的玻璃管下端有一个球形小容器，管内有一段长约 $1\Ucm$ 的水银柱。现在需要用比较准确的方法测定球形小容器的容积 $V$。可用的器材有：刻度尺（量程 $500\Umm$）、温度计（测量范围 $0\sim100\UdC$）、玻璃容器（高约为 $30\Ucm$，直径约 $10\Ucm$）、足够多的沸水和冷水。
\begin{enumerate}
	\item
	简要写出实验步骤及需要测量的物理量；
	\item
	说明如何根据所测得的物理量得出实验结果。
\end{enumerate}
\onepicture[h=5cm, align=right]{\includesvg{figs/17}}

%% answer:
\jdanswer{
\begin{enumerate}
	\item
	将水银柱以下的玻璃管浸没在水中，改变水温，用温度计测得若干组（或两组）不同水温（即气体温度）$T$ 和气体长度 $x$ 的值。
	\item
	方法一：气体作等压变化，有 $(V+xS)=CT$，即 $xS=CT-V$；作 $xS-T$ 图，图像截距的绝对值即为 $V$。方法二：测两组数据，有 $x_{1}S=CT_{1}-V$，$x_{2}S=CT_{2}-V$；得 $V=\frac{T_{1}x_{2}-T_{2}x_{1}}{T_{2}-T_{1}}S$。
\end{enumerate}
}

%% memo:
\memoanswer{
（1）将水银柱以下的玻璃管浸没在水中，改变水温，用温度计测得若干组（或两组）不同水温（即气体温度）$T$ 和气体长度 $x$ 的值。

（2）方法一：气体作等压变化，有
\[
(V+xS)=CT,
\]
即
\[
xS=CT-V,
\]
作 $xS-T$ 图，图像截距的绝对值即为 $V$。

方法二：测两组数据，有
\[
x_{1}S=CT_{1}-V,\qquad x_{2}S=CT_{2}-V,
\]
得
\[
V=\frac{T_{1}x_{2}-T_{2}x_{1}}{T_{2}-T_{1}}S.
\]
}

\item
%% number: 18
%% paperName: 2004年上海高考第 18 题 10 分
%% typeId: 4
%% type: 实验题
%% chapter: 电路
%% point: 分压与分流
%% method: 图象法
%% score: 10
%% degree: 450
%% duplicateId: 0
%% body:
小灯泡灯丝的电阻会随温度的升高而变大。某同学为研究这一现象，用实验得到如下数据（$I$ 和 $U$ 分别表示小灯泡上的电流和电压）：

\begin{center}
\begin{tabular}{|c|c|c|c|c|c|c|c|c|c|c|}
\hline
$I$（A） & 0.12 & 0.21 & 0.29 & 0.34 & 0.38 & 0.42 & 0.45 & 0.47 & 0.49 & 0.50 \\
\hline
$U$（V） & 0.20 & 0.40 & 0.60 & 0.80 & 1.00 & 1.20 & 1.40 & 1.60 & 1.80 & 2.00 \\
\hline
\end{tabular}
\end{center}

\twopicture[showlabel=false, wA=4.5cm, wB=4.5cm]{figs/18a.png}{\includesvg{figs/18b}}
\begin{enumerate}
	\item
	在左下框中画出实验电路图。可用的器材有：电压表、电流表、滑线变阻器（变化范围 $0\sim10\UO$）、电源、小灯泡、电键、导线若干。
	\item
	在右图中画出小灯泡的 $U-I$ 曲线。
	\item
	如果实验中测得电池的电动势是 $1.5\UV$，内阻是 $2.0\UO$，问：将本题中的小灯泡接在该电池的两端，小灯泡的实际功率是多少？（简要写出求解过程：若需作图，可直接画在第（2）小题方格图中）
\end{enumerate}

%% answer:
\jdanswer{
\begin{enumerate}
	\item
	分压器接法。
	\item
	如图。
	\item
	作出 $U=E-Ir$ 图线，可得小灯泡工作电流为 $0.35\UA$，工作电压为 $0.80\UV$，实际功率为 $0.28\UW$。
\end{enumerate}
}

%% memo:
\memoanswer{
（1）实验电路图如图所示。

\onepicture[w=4.5cm]{\includesvg{figs/18_an1}}

（2）$U-I$ 曲线如图所示。

\onepicture[w=6cm]{\includesvg{figs/18_an2}}

（3）设灯泡两端的电压为 $U$，通过的电流为 $I$，由闭合电路的欧姆定律知，$U=E-Ir$，即 $U=1.5-2I$；作出 $U-I$ 图线，由图线的交点可得小灯泡工作电流为 $0.35\UA$，工作电压为 $0.80\UV$，因此小灯泡实际功率为 $0.28\UW$。
}

\item
%% number: 19
%% paperName: 2004年上海高考第 19 题 8 分
%% typeId: 5
%% type: 辨析题
%% chapter: 电场
%% point: 库仑定律
%% method: 无
%% score: 8
%% degree: 650
%% duplicateId: 0
%% body:
“真空中两个整体上点电荷相距 $10\Ucm$，它们之间相互作用力大小为 $9\times10^{-4}\UN$。当它们合在一起时，成为一个带电量为 $3\times10^{-8}\UC$ 的点电荷。问原来两电荷的带电量各是多少？”

某同学求解如下：

根据电荷守恒定律：$q_{1}+q_{2}=3\times10^{-8}\UC=a$　　（1）

根据库仑定律：$q_{1}q_{2}=\frac{Fr^{2}}{k}=\frac{9\times10^{-4}\times(10\times10^{-2})^{2}}{9\times10^{9}}\UC^{2}=1\times10^{-15}\UC^{2}=b$

以 $q_{2}=\frac{b}{q_{1}}$ 代入（1）式得：$q_{1}^{2}-aq_{1}-b=0$

解得 $q_{1}=\frac{1}{2}\left(a\pm\sqrt{a^{2}-4b}\right)=\frac{1}{2}\left(3\times10^{-8}\pm\sqrt{9\times10^{-16}-4\times10^{-15}}\right)\UC$

根号中的数值小于 $0$，经检查，运算无误。试指出求解过程中的问题并给出正确的解答。

%% answer:
\jdanswer{
题中仅给出相互作用力的大小，两点电荷可能异号，按电荷异号计算。$q_{1}=5\times10^{-8}\UC$，$q_{2}=2\times10^{-8}\UC$。
}

%% memo:
\memoanswer{
题中仅给出相互作用力的大小，两点电荷可能异号，按电荷异号计算。

$q_{1}-q_{2}=3\times10^{-8}\UC=a$

$q_{1}q_{2}=\frac{Fr^{2}}{k}=\frac{9\times10^{-4}\times(10\times10^{-2})^{2}}{9\times10^{9}}\UC^{2}=1\times10^{-15}\UC^{2}=b$

以 $q_{2}=\frac{b}{q_{1}}$ 代入（1）式得：$q_{1}^{2}-aq_{1}-b=0$

解得：$q_{1}=5\times10^{-8}\UC$，$q_{2}=2\times10^{-8}\UC$。
}

\item
%% number: 20
%% paperName: 2004年上海高考第 20 题 12 分
%% typeId: 6
%% type: 计算题
%% chapter: 气体性质
%% point: 玻意耳定律
%% method: 无
%% score: 12
%% degree: 450
%% duplicateId: 0
%% body:
如图所示，一端封闭、粗细均匀的薄壁玻璃管开口向下竖直插在装有水银的水银槽内，管内封闭有一定质量的空气，水银槽的截面积上下相同，是玻璃管截面积的 5 倍。开始时管内空气长度为 $6\Ucm$，管内外水银面高度差为 $50\Ucm$。将玻璃管沿竖直方向缓慢上移（管口未离开槽中水银），使管内外水银面高度差变成 $60\Ucm$。（大气压相当于 $75\UcmHg$），求：
\begin{enumerate}
	\item
	此时管内空气柱的长度；
	\item
	水银槽内水银面下降的高度。
\end{enumerate}
\onepicture[h=4.5cm, align=right]{\includesvg{figs/20}}

%% answer:
\jdanswer{
\begin{enumerate}
	\item
	$l_{2}=0.10\Um$。
	\item
	$\Delta x=0.02\Um$。
\end{enumerate}
}

%% memo:
\memoanswer{
（1）管内空气作等温变化，由玻意耳定律得
\[
(p_{0}-\rho gH_{1})l_{1}=(p_{0}-\rho gH_{2})l_{2},
\]
\[
l_{2}=\frac{p_{0}-\rho gH_{1}}{p_{0}-\rho gH_{2}}l_{1}=\frac{75-50}{75-60}\times0.06\Um=0.10\Um.
\]

（2）设槽内水银面下降 $\Delta x$，由水银体积不变得
\[
S_{1}\Delta H=(S_{2}-S_{1})\Delta x,\qquad \Delta H+\Delta x=H,
\]
解得
\[
\Delta x=0.2H=0.02\Um.
\]
}

\item
%% number: 21
%% paperName: 2004年上海高考第 21 题 12 分
%% typeId: 6
%% type: 计算题
%% chapter: 曲线运动
%% point: 平抛运动
%% method: 分情况讨论
%% score: 12
%% degree: 450
%% duplicateId: 0
%% body:
滑雪者从 A 点由静止沿斜面滑下，经一平台水平飞离 B 点，地面上紧靠着平台有一个水平台阶，空间几何尺度如图所示。斜面、平台与滑雪板之间的动摩擦因数为 $\mu$，假设滑雪者由斜面底端进入平台后立即沿水平方向运动，且速度大小不变。求：
\begin{enumerate}
	\item
	滑雪者离开 B 点时的速度大小；
	\item
	滑雪者从 B 点开始做平抛运动的水平距离 $s$。
\end{enumerate}
\onepicture[w=7cm, align=right]{\includesvg{figs/21}}

%% answer:
\jdanswer{
\begin{enumerate}
	\item
	$v_{B}=\sqrt{2g(H-h-\mu L)}$。
	\item
	$s_{1}=\sqrt{2h(H-h-\mu L)}$，此时必须满足 $H-\mu L<2h$；当 $H-\mu L>2h$ 时，滑雪者直接落到地面上，$s_{2}=2\sqrt{h(H-h-\mu L)}$。
\end{enumerate}
}

%% memo:
\memoanswer{
（1）设滑雪者质量为 $m$，斜面与水平面夹角为 $\theta$，斜面长为 $s$，滑雪者在滑行过程中克服摩擦力做的功为
\[
W_{f}=\mu mgs\cos\theta+\mu mg(L-s\cos\theta)=\mu mgL,
\]
由动能定理
\[
mg(H-h)-\mu mgL=\frac{1}{2}mv_{B}^{2},
\]
得
\[
v_{B}=\sqrt{2g(H-h-\mu L)}.
\]

（2）设滑雪者离开 B 点后落在台阶上，由平抛规律得
\[
\frac{h}{2}=\frac{1}{2}gt_{1}^{2},\qquad s_{1}=v_{B}t_{1}<\sqrt{2}h,
\]
得
\[
s_{1}=\sqrt{2h(H-h-\mu L)},
\]
此时必须满足 $H-\mu L<2h$。

当 $H-\mu L>2h$ 时，滑雪者直接落到地面上，由平抛规律得
\[
h=\frac{1}{2}gt_{2}^{2},\qquad s_{2}=v_{B}t_{2},
\]
得
\[
s_{2}=2\sqrt{h(H-h-\mu L)}.
\]
}

\item
%% number: 22
%% paperName: 2004年上海高考第 22 题 14 分
%% typeId: 6
%% type: 计算题
%% chapter: 电磁感应
%% point: 电磁感应中的变加速
%% method: 图象法
%% score: 14
%% degree: 450
%% duplicateId: 0
%% body:
水平面上两根足够长的金属导轨平行固定放置，间距为 $L$，一端通过导线与阻值为 $R$ 的电阻连接；导轨上放一质量为 $m$ 的金属杆（见右上图），金属杆与导轨的电阻不计；均匀磁场竖直向下。用与导轨平行的恒定力 $F$ 作用在金属杆上，杆最终将做匀速运动。当改变拉力的大小时，相对应的匀速运动速度 $v$ 也会改变，$v$ 和 $F$ 的关系如右下图。（取重力加速度 $g=10\Umsq$）
\begin{enumerate}
	\item
	金属杆在匀速运动之前做作什么运动？
	\item
	若 $m=0.5\Ukg$，$L=0.5\Um$，$R=0.5\UO$，磁感应强度 $B$ 为多大？
	\item
	由 $v-F$ 图线的截距可求得什么物理量？其值为多少？
\end{enumerate}
\onepicture[h=2cm, align=right]{\includesvg{figs/22a}}
\onepicture[h=5cm, align=right]{\includesvg{figs/22b}}

%% answer:
\jdanswer{
\begin{enumerate}
	\item
	金属杆做变速运动（或变加速运动、加速度减小的加速运动）。
	\item
	$B=1\UT$。
	\item
	由截距可以求得杆所受阻力为 $f=2\UN$，动摩擦因数为 $0.4$。
\end{enumerate}
}

%% memo:
\memoanswer{
（1）金属杆做变速运动（或变加速运动、加速度减小的加速运动）。

（2）金属杆切割产生的感应电动势为
\[
E=BLv,
\]
感应电流为
\[
I=\frac{E}{R},
\]
金属杆受到的安培力为
\[
F_{A}=BIL=\frac{B^{2}L^{2}v}{R},
\]
由图线可知杆受拉力、安培力和阻力作用，匀速时合力为零，有
\[
F=\frac{B^{2}L^{2}v}{R}+f,
\]
所以
\[
v=\frac{R}{B^{2}L^{2}}F-\frac{fR}{B^{2}L^{2}},
\]
由图线直接得出斜率为 $k=2$，联立解得
\[
B=\sqrt{\frac{R}{kL^{2}}}=1\UT.
\]

（3）由直线的截距可以求得金属杆受到的阻力 $f$，由纵截距知 $f=2\UN$。若金属杆受到的阻力仅为滑动摩擦力，由截距可求得动摩擦因数 $\mu=0.4$。
}

\item
%% number: 23
%% paperName: 2004年上海高考第 23 题 14 分
%% typeId: 6
%% type: 计算题
%% chapter: 力物体的平衡
%% point: 力矩平衡
%% method: 无
%% score: 14
%% degree: 250
%% duplicateId: 0
%% body:
有人设计了一种新型伸缩拉杆秤。结构如图，秤杆的一端固定一配重物并悬一挂钩，秤杆外面套有内外两个套筒，套筒左端开槽使其可以不受秤纽阻碍而移动到挂钩所在的位置（设开槽后套筒的重心仍在其长度中点位置）。秤杆与内层套筒上刻有质量刻度。空载（挂钩上不挂物体，且套筒未拉出）时，用手提起秤纽，杆杆秤恰好平衡。当物体挂在挂钩上时，往外移动内外套筒可使杆秤平衡，从内外套筒左端的位置可以读得两个读数，将这两个读数相加，即可得到待测物体的质量。已知秤杆和两个套筒的长度均为 $16\Ucm$，套筒可移出的最在距离为 $15\Ucm$，秤纽到挂钩的距离为 $2\Ucm$，两个套筒的质量均为 $0.1\Ukg$。取重力加速度 $g=10\Umsq$。求：
\begin{enumerate}
	\item
	当杆秤空载平衡时，秤杆、配重物及挂钩所受重力相对秤纽的合力矩；
	\item
	当在秤钩上挂一物体时，将内套筒向右移动 $5\Ucm$，外套筒相对内套筒向右移动 $8\Ucm$，杆秤达到平衡，物体的质量多大？
	\item
	若外层套筒不慎丢失，在称某一物体时，内层套筒的左端在读数为 $1\Ukg$ 处杆秤恰好平衡，则该物体实际质量多大？
\end{enumerate}
\twopicture[showlabel=false, wA=7cm, wB=7cm]{figs/23a.png}{figs/23b.png}

%% answer:
\jdanswer{
\begin{enumerate}
	\item
	$M=0.12\UNcdotm$。
	\item
	$m_{1}=0.9\Ukg$。
	\item
	$m_{2}=0.2\Ukg$。
\end{enumerate}
}

%% memo:
\memoanswer{
（1）套筒不拉出时杆恰平衡，两套筒相对秤纽的力矩与所求力矩相等，设套筒长为 $L$，合力矩为
\[
M=2mg\left(\frac{L}{2}-d\right)=2\times0.1\times10\times(0.08-0.02)\UNcdotm=0.12\UNcdotm.
\]

（2）由力矩平衡得
\[
m_{1}gd=mgx_{1}+mg(x_{1}+x_{2}),
\]
解得
\[
m_{1}=\frac{2x_{1}+x_{2}}{d}m=\frac{2\times0.05+0.08}{0.02}\times0.1\Ukg=0.9\Ukg.
\]

（3）正常称 $1\Ukg$ 物体时内外套筒可一起向外拉出 $x'$，由力矩平衡得
\[
m_{2}'gd=2mgx',
\]
解得
\[
x'=\frac{m_{2}'}{2m}d=\frac{1}{2\times0.1}\times0.02\Um=0.1\Um.
\]
外层套筒丢失后称物，此时内套筒左端离秤纽距离为
\[
x'-d=0.08\Um,
\]
由力矩平衡得
\[
m_{2}gd+M=mg\left(x'-d+\frac{L}{2}\right),
\]
解得
\[
m_{2}=\frac{m}{d}\left(x'-d+\frac{L}{2}\right)-\frac{M}{gd}=0.2\Ukg.
\]
}

\end{enumerate}

\end{document}
